\originalpreface
电动力学把电荷和电流产生的场、场对物质的作用及能量与动量的传递放在同一个体系中. 本稿面向已经学过普通物理电磁学、矢量微积分和常微分方程的读者, 以本科边值问题、麦克斯韦方程和辐射为学习终点. 所用原件来自 MIT OpenCourseWare, 中文陈述、推导组织和补充解答为本仓库重编; 本稿并非 MIT 官方中文讲义, 也不表示 MIT 认可.

\noindent\textbf{主课程.} 8.07 Electromagnetism II, Fall 2012, 教师 Alan Guth 与 Min Chen. 官方先修为 8.03 Vibrations and Waves 和 18.03 Differential Equations; 每周三次一小时讲课. 教材为 David J. Griffiths, \textit{Introduction to Electrodynamics}, 第三版. 教材本身不在原件包中. 考核为每周一次作业, 共十次; 两次 Quiz 和课堂期末考试. 作业占 25\%, Quiz 合计占 40\%, 期末占 35\%. 允许讨论作业, 但须独立写出解答.

\noindent\textbf{课历与大纲.} 官方课历只给讲次及相对作业节点, 本表不添加未公布的日期.
\begin{longtable}{p{.08\textwidth}p{.60\textwidth}p{.23\textwidth}}
\toprule 讲次 & 官方主题的中文表述 & 作业与考试节点\\\midrule
1--3 & 矢量分析; 指标记号中的转动; 基本方程的积分形式 & 第2讲布置作业1\\
4--7 & 导体表面力与电容; 镜像法; 分离变量; 球坐标分离变量; 带孔球壳 & 布置作业2--4; 作业1、2到期\\
8--10 & 无迹对称张量; 勒让德多项式与球谐函数; 磁多极与磁偶极 & 布置作业5、6; 作业3--5到期; Quiz 1\\
11--13 & 物质中的磁场; 超材料的反向切伦科夫辐射 & 布置作业7--9; 作业6、7到期\\
14--17 & 电磁波; 势与场; 辐射; 波、势及辐射复习 & 布置作业10; Quiz 2; 作业8--10到期; Final\\\bottomrule
\end{longtable}

\clearpage
\noindent\textbf{进阶课程.} 8.311 Electromagnetic Theory, Spring 2004, 教师 Leonid Levitov, 属研究生课程. 每周两次一小时半讲课及一次一小时半习题课, 先修为 8.07. 十二次作业占 60\%, 闭卷期末占 40\%. 作业官方解答由助教 Chiyan Luo 提供（文件7内容错配为作业8, 第7套匹配解未取得）, 保留署名 \textit{Courtesy of Chiyan Luo. Used with permission.} 课程使用商业教材, OCW 只提供四份补充笔记, 不能把这四份标成全部课堂理论原稿.

其官方课历为: L1--4 麦克斯韦方程、能量与动量守恒、电磁波的能量动量; L5--8 电响应模型、色散、Kramers--Kronig 关系及菲涅耳公式; L9--10 传输线、波导与谐振腔; L11--13 粒子与场的拉格朗日、哈密顿及 Noether 定理; L14--15 洛伦兹变换; L16--20 推迟势、偶极辐射、散射、相对论辐射与辐射反作用; L21--25 辐射频谱、相对论拉莫尔公式、脉冲辐射、切伦科夫辐射和同步辐射; 最后为期末考试. R1--R12 为习题课. 本稿进阶部分与本科部分明确分列.

\noindent\textbf{来源与解答身份.} 已校验的原始资料为 61 份 PDF, 共 474 页; 另两份原件只保留链接、哈希及页数, 不打包整件. 8.07 的两次 Quiz 有官方解答, 本稿对应部分为中文整理并补充步骤. 研究生12套的66道一级题已逐题纳入后文, 解答为AI独立推导并参照官方解答核对; 官方缺解与确定的勘误在题旁注明. 十份本科作业与本科期末未取得官方解答, 相应解答标为 AI 独立解答. 商业教材题号本身不是题面; 仅给题号且缺少完整条件的项目保持缺失状态. 本文采用 SI 制, $c^{-2}=\eps\mu_0$, 除明确介质章节外 $\ee,\bb$ 指真空场. 静态理想偶极的接触项按球形主值理解.

原件按 MIT OCW 条款与 CC BY-NC-SA 4.0 使用, 第三方例外不并入默认许可. 第13讲特别保留 \textit{Courtesy of Hongsheng Chen and Min Chen. Used with permission.} 讲义中所有新图为根据明确模型用 TikZ 绘制的示意图, 不恢复 OCW 已移除的图片.
